\section{Информация о курсе} 
\subsection{Оценка} 

Формула итоговой оценки за курс:
\[ 
\begin{aligned}
\text{О} ={} & 0.15 \cdot \text{ДЗ} + 0.1 \cdot \text{Quiz} + 0.15 \cdot \text{Контест} + {} \\
& 0.2 \cdot \text{КР (БЛОК)} + 0.25 \cdot \text{Коллок (БЛОК)} + {} \\
& 0.05 \cdot \text{Бонус от семинариста} + 0.2 \cdot \text{Коллок по доп.\ материалу} 
\end{aligned}
\]

\section{Введение}

\begin{definition}[Случайный эксперимент]
    \textbf{Случайный эксперимент} --- это процесс, наблюдение или действие, результат которого невозможно точно предсказать заранее (до его проведения). 
    
    Чтобы эксперимент можно было изучать методами теории вероятностей, он должен обладать следующими свойствами:
    \begin{itemize}
        \item \textbf{Отсутствие детерминированной регулярности:} при многократном повторении эксперимента в неизменных условиях результаты могут различаться (исход единичного опыта случаен).
        \item \textbf{Воспроизводимость эксперимента:} эксперимент теоретически или практически можно повторить бесконечное число раз в одних и тех же контролируемых условиях. Его проведение и исход не должны быть завязаны на уникальные свойства или волю конкретной личности.
        \item \textbf{Статистическая устойчивость:} при очень большом числе повторений частота появления того или иного исхода стабилизируется и стремится к фиксированному числу (вероятности).
        \[
        \frac{N_{1, N}}{N_1} \approx \frac{N_{2, N}}{N_2} \approx \ldots \approx \frac{N_{k, N}}{N_k}
        \]
    \end{itemize}
\end{definition}

Сопоставим каждому результату эксперимента свою $\omega$ --- \textbf{элементарный исход}. Множество всех элементарных исходов обозначается как $\Omega$:
\[
    \Omega = \{\omega\} 
\]

\begin{example}[Пример: Бросок двух игральных костей]
    Кидаются 2 кости, нас интересует сумма. Варианты выбора пространства $\Omega$:
    \begin{enumerate}[label=\arabic*)]
        \item $\omega \in \{2, \dots, 12\}, \quad |\Omega| = 11$
        \item $\omega = (i_1, i_2), \quad i_j \in \{1, \dots, 6\}, \quad |\Omega| = 36$
        \item $\omega = \{i_1, i_2\}, \quad |\Omega| = \frac{36-6}{2} + 6 = 21$
    \end{enumerate}
\end{example}

\subsection{События}

\begin{definition}[Случайное событие]
    \textbf{Случайное событие} --- это подмножество пространства элементарных исходов: $\mathcal{A} \subset \Omega$.
\end{definition}

\begin{definition}[Вероятность]
    \textbf{Вероятность} события $\mathcal{A}$ --- это предел относительной частоты его появления при стремлении числа экспериментов к бесконечности:
    \[
        \frac{N_{\mathcal{A}}}{N} \xrightarrow[N \to \infty]{} \Pr(\mathcal{A})
    \]
\end{definition}

\subsection{Математические модели}

Любая математическая модель в теории вероятностей состоит из трех объектов: $(\Omega, \mathcal{F}, \Pr)$, где: 
\begin{itemize}
    \item $\Omega$ --- пространство элементарных исходов;
    \item $\mathcal{F}$ --- множество событий, такое что $\mathcal{F} \subset 2^{\Omega}$ ($\sigma$-алгебра);
    \item $\Pr$ --- вероятностная мера, задающая функцию $\Pr \colon \mathcal{F} \to [0; 1]$.
\end{itemize}

\begin{exercise}[О выборе $\mathcal{F}$]
    Всегда ли $\mathcal{F} = 2^{\Omega}$? Верно ли, что $\Omega \in \mathcal{F}$?
\end{exercise}

\subsubsection{Свойства вероятностной меры}

\begin{itemize}
    \item Нормированность: $\Pr(\Omega) = 1$.
    \item Вероятность невозможного события: $\Pr(\varnothing) = 0$. 
    \item \textit{(Отсылка к линейности: $\Pr(\varnothing) = 0 \implies f(\alpha \vec{x} + \beta \vec{y}) = \alpha f(\vec{x}) + \beta f(\vec{y})$)}.
    \item \textbf{Аддитивность:} для несовместных событий (т.е. $\mathcal{A} \cap \mathcal{B} = \varnothing$):
    \[
        \Pr(\mathcal{A} \sqcup \mathcal{B}) = \Pr(\mathcal{A}) + \Pr(\mathcal{B}),
    \]
    где $\sqcup$ --- дизъюнктивное объединение.
\end{itemize}

\subsection{Дискретные модели}

Пусть $\Omega$ --- конечное множество. Тогда:
\[
    \mathcal{F} = 2^{\Omega} \implies \forall \omega \in \Omega \colon \{\omega\} \in \mathcal{F}, \quad |\mathcal{F}| = 2^{|\Omega|}
\]

Как задается $\Pr$? Достаточно задать вероятности каждого исхода:
\[
    \forall \omega \in \Omega \colon \Pr(\omega)
\]
Тогда для любого события $\mathcal{A} \in \mathcal{F}$ ($\mathcal{A} \subset \Omega$) вероятность равна сумме вероятностей входящих в него исходов:
\[
    \Pr(\mathcal{A}) = \Sum{\omega \in \mathcal{A}}{} \Pr(\omega)
\]

\subsubsection{Классические дискретные модели}

\begin{definition}[Классическая модель]
    \textbf{Классическая модель} --- это дискретная модель, в которой все исходы равновероятны: $\Pr(\omega) = \text{const}$.
\end{definition}

Найдем эту константу из свойства нормированности:
\[
    1 = \Pr(\Omega) = \Pr\l(\bigsqcup^{|\Omega|}_{i = 1} \{\omega_i\}\r) = \Sum{i = 1}{|\Omega|} \Pr(\omega_i) = |\Omega| \cdot \text{const} \implies \text{const} = \frac{1}{|\Omega|}
\]
Следовательно, вероятность любого события в классической модели:
\[
    \Pr(\mathcal{A}) = \Sum{\omega \in \mathcal{A}}{} \Pr(\omega) = \Sum{\omega \in \mathcal{A}}{} \frac{1}{|\Omega|} = \frac{|\mathcal{A}|}{|\Omega|}
\]

\begin{example}[Пример: Два кубика в классической модели]
    \begin{itemize}
        \item Различимые кубики: $\omega = (i_1, i_2), \quad |\Omega| = 36$.
        \[
            \mathcal{F} = 2^{\Omega}, \quad \Pr(\mathcal{A}) = \frac{|\mathcal{A}|}{36}
        \]
        \item Неразличимые кубики: $\omega = \{i_1, i_2\}, \quad |\Omega| = 21$.
        \[
            \mathcal{F} = 2^{\Omega}, \quad \Pr(\omega) \neq \text{const} \quad \text{(модель перестает быть классической!)}
        \]
    \end{itemize}
\end{example}

\begin{task}[Урновая схема: выбор с возвращением]
    Есть $N$ шариков: $M$ --- красных (1), $N - M$ --- черных (0). 
    Вытаскиваем $n$ шариков с возвращением. Какова вероятность вытащить ровно $k$ красных шариков (в любом порядке)?
\end{task}

\begin{solution}
    Если мы обозначим элементарный исход как: 
    \[
        \omega = (i_1, \dots, i_n), \quad i_j \in \{0, 1\}
    \]
    то исходы перестанут быть равновероятными. 

    Рассмотрим элементарный исход как упорядоченный набор вытащенных шариков:
    \[
        \omega = (i_1, \dots, i_n), \quad i_j \in \{1, \dots, N\}
    \]
    где номера $1, \dots, M$ --- красные, а $M+1, \dots, N$ --- черные шары.
    
    Общее число исходов:
    \[
        |\Omega| = N^n
    \]
    Число благоприятных исходов (выбираем $k$ позиций для красных шаров из $n$ возможных):
    \[
        |\mathcal{A}| = C_n^k \cdot M^k \cdot (N - M)^{n - k}
    \]
    Вероятность события:
    \[
        \Pr(\mathcal{A}) = \frac{|\mathcal{A}|}{|\Omega|} = \frac{C_n^k \cdot M^k \cdot (N - M)^{n-k}}{N^n} = C_n^k \l( \frac{M}{N} \r)^k \l(1 - \frac{M}{N} \r)^{n-k}  
    \]
    Если обозначить $p = \frac{M}{N}$, получим классическую формулу Бернулли:
    \[
        \Pr(\mathcal{A}) = C_n^k \, p^k (1 - p)^{n-k} 
    \]
\end{solution}

\subsubsection{Неклассические дискретные модели}

\begin{definition}[Схема испытаний Бернулли]
    Обозначается $\mathrm{Bern}(n, p)$, где $n \in \N$ --- число испытаний, $p \in (0; 1)$ --- вероятность успеха. Исход каждого испытания: 0 (неудача) или 1 (успех).
\end{definition}

Элементарный исход --- вектор результатов:
\[
    \omega = (i_1, \dots , i_n), \quad i_j \in \{0, 1\} \implies |\Omega| = 2^n
\]
Вероятность конкретного исхода $\omega$ (исходы не равновероятны):
\[
    \Pr(\omega) = p^{\Sum{j=1}{n} i_j} \cdot (1 - p)^{n - \Sum{j=1}{n} i_j}
\]

Проверим нормированность (сумма вероятностей всех исходов равна 1):
\[
\begin{aligned}
    \Sum{\omega \in \Omega}{} \Pr(\omega) &= \Sum{\omega \in \Omega}{} p^{\Sum{j=1}{n} i_j} \cdot (1 - p)^{n - \Sum{j=1}{n} i_j} \\
    &= \Sum{k = 0}{n} \left( \Sum{\omega \colon \Sum{j = 1}{n} i_j = k}{} p^k (1-p)^{n-k} \right) \\
    &= \Sum{k = 0}{n} C_n^k \, p^k (1-p)^{n-k} \\
    &= (p + (1 - p))^n = 1^n = 1
\end{aligned}
\]

\subsection{Недискретные модели}
\subsubsection{Геометрическая вероятность}

Если $\Omega \subset \R^m$ имеет бесконечное (несчетное) число исходов, то вероятности задаются через меру (длину, площадь, объем).

\begin{example}[Геометрическая вероятность на отрезке]
    Пусть $\Omega = [0; 1]$.
    Из-за наличия неизмеримых множеств $\mathcal{F} \neq 2^{\Omega}$, поэтому вероятность попадания в измеримое множество $\mathcal{A} \subset \Omega$ определяется как:
    \[
        \Pr(\mathcal{A}) = \frac{\lambda(\mathcal{A})}{\lambda(\Omega)} = \lambda(\mathcal{A}),
    \]
    где $\lambda(\mathcal{A})$ --- мера Лебега множества $\mathcal{A}$.
\end{example}
